% !TeX root = ../../main.tex
\subsection{\texorpdfstring{BARCS generalizes the CB\textsuperscript{2}
sampling model}{BARCS generalizes the CB2 sampling model}}
\label{sec:cb2-generalization}

CB\textsuperscript{2} established that a guide count is naturally analyzed as
a proportion conditional on the corresponding library total.  This
conditioning matters because the beta-binomial variance separates two sources
of uncertainty: binomial sampling variation, which decreases with sequencing
depth, and between-library heterogeneity, which does not.  A
negative-binomial model can accommodate overdispersion in marginal counts, but
it does not encode the observed library total as the denominator of the guide
proportion.  Prior goodness-of-fit analyses therefore support the
beta-binomial sampling model for pooled-guide counts
\citep{jeong2019beta}.  The present work asks a different question: can that
sampling model support regression designs beyond the two groups handled by
CB\textsuperscript{2}?

BARCS changes the mean model while retaining the sampling model.  For guide
$g$ in library $i$,
\[
  K_{gi}\mid N_i \sim
  \operatorname{BetaBinomial}(N_i,\mu_{gi},\rho_g),
  \qquad
  \logit(\mu_{gi})=\mathbf{x}_i^\top\boldsymbol{\beta}_g.
\]
The original two-group question is recovered when $\mathbf{x}_i$ contains an
intercept and one group indicator.  More generally, the columns of
$\mathbf{X}$ can represent time, dose, donor, batch, an interaction, or a
prespecified contrast.  Thus BARCS is a generalization of the statistical
question that CB\textsuperscript{2} can ask; it is not claimed to outperform
CB\textsuperscript{2} on the binary design for which the original method was
developed.  Supplementary Methods establish exact equivalence on the legacy
group-summary scale and first-order equivalence of the default logit statistic
under local alternatives.
